\section*{الفصل \ref{ch:b1:acetals-protection} --- مركبات الكربونيل: الإضافات المحبة للنواة، والأسيتالات والحماية}

\begin{solution}{exo:b1:acetals-protection:1}
\ce{CH3CH(OH)OCH3}: \omterm{def:b1:acetals-protection:hemiacetal}{نصف أسيتال}. \ce{CH3CH(OCH3)2}: \omterm{def:b1:acetals-protection:hemiacetal}{أسيتال}.
\ce{CH3CH(OH)2}: هيدرات. \ce{CH3OCH3}: إيثر. \ce{(CH3)2C(OCH2CH3)2}:
\omterm{def:b1:acetals-protection:hemiacetal}{أسيتال} (لكيتون).
\end{solution}

\begin{solution}{exo:b1:acetals-protection:2}
\ce{CH3CH2CHO + 2CH3OH <=> CH3CH2CH(OCH3)2 + H2O}؛ والبروبانون مع
الإيثان-1,2-ثنائي أول يعطي 2,2-ثنائي ميثيل-1,3-ديوكسولان والماء،
\ce{(CH3)2CO + HOCH2CH2OH <=> C5H10O2 + H2O}.
\end{solution}

\begin{solution}{exo:b1:acetals-protection:3}
هيدروكسيد الصوديوم، وبروميد ميثيل المغنيزيوم وبوروهيدريد الصوديوم تتركه
دون تغيير؛ أما حمض الكبريتيك المائي المخفف فيحلمئه.
\end{solution}

\begin{solution}{exo:b1:acetals-protection:4}
تُضاف OH الكربون 4 إلى كربون الألدهيد 1: حلقة خماسية من
أربع ذرات كربون وذرة أكسجين واحدة، مع مجموعة OH على الكربون المجاور لأكسجين
الحلقة.
\end{solution}

\begin{solution}{exo:b1:acetals-protection:5}
برتنة C=O؛ ثم يُضاف الميثانول إلى الكربون؛ ثم فقد \ce{H+}: \omterm{def:b1:acetals-protection:hemiacetal}{نصف أسيتال}؛
ثم برتنة OH فيه؛ ثم فقد الماء نحو \ce{CH3CH=O+CH3}؛ ثم يُضاف ميثانول
ثانٍ؛ ثم فقد \ce{H+}: \ce{CH3CH(OCH3)2}. والبروتونان المأخوذان يُعادان:
فالحمض \omterm{def:b1:elementary-steps:catalyst}{حفّاز}.
\end{solution}

\begin{solution}{exo:b1:acetals-protection:6}
برتنة إحدى مجموعتي OCH3؛ ثم فقد الميثانول نحو \ce{(CH3)2C=O+CH3}؛ ثم يُضاف الماء؛
ثم فقد \ce{H+}: \omterm{def:b1:acetals-protection:hemiacetal}{نصف أسيتال}؛ ثم برتنة، وفقد الميثانول، وفقد \ce{H+}:
البروبانون. والفائض من الماء يُبقي $Q$ أدنى من $K$ للحلمأة، التي تمضي
حتى الاكتمال.
\end{solution}

\begin{solution}{exo:b1:acetals-protection:7}
جزيء ماء واحد لكل \omterm{def:b1:acetals-protection:hemiacetal}{أسيتال}: $0.250 \times 18.0 = \qty{4.5}{g}$، أي نحو
\qty{4.5}{mL}. وبعد ساعة، $3.2/18.0 = \qty{0.178}{mol}$: تحول بنسبة
\qty{71}{\percent}.
\end{solution}

\begin{solution}{exo:b1:acetals-protection:8}
مع ثنائي الأول، يحل جزيء واحد محل جزيئين: فالتفاعل لا يُنقص
عدد الجزيئات (ألدهيد واحد وثنائي أول واحد يعطيان أسيتالًا واحدًا وجزيء
ماء واحدًا)، ومجموعة OH الثانية مثبتة قريبًا من الكربون المتفاعل، فتنغلق
الحلقة بسهولة. وكلا الأمرين يجعل \omterm{def:b1:acetals-protection:hemiacetal}{الأسيتال} الحلقي أكثر ملاءمة.
\end{solution}

\begin{solution}{exo:b1:acetals-protection:9}
كان الحمض سيهدم \omterm{def:b1:organomagnesium:organometallic}{كاشف غرينيار} (انتقال بروتون)؛ ويجب أيضًا إزالة آثار
الماء وثنائي الأول.
\end{solution}

\begin{solution}{exo:b1:acetals-protection:10}
1. احمِ كيتون 4-كلوروبوتان-2-ون على هيئة \omterm{def:b1:acetals-protection:hemiacetal}{أسيتال} حلقي
(الإيثان-1,2-ثنائي أول، \omterm{def:b1:elementary-steps:catalyst}{حفّاز} حمضي، دين--ستارك). 2. المغنيزيوم في إيثر جاف:
\omterm{def:b1:organomagnesium:organometallic}{كاشف غرينيار} للكلوريد المحمي، الذي لم يعد كربونيله
الخاص يهدمه. 3. أضف البروبانون؛ ثم المعالجة بكلوريد الأمونيوم المائي:
الكحول الثالثي. 4. حمض مائي مخفف: \omterm{def:b1:acetals-protection:protecting-group}{إزالة الحماية} إلى
5-هيدروكسي-5-ميثيل هكسان-2-ون.
\end{solution}

\begin{solution}{exo:b1:acetals-protection:11}
$Q = [\text{أسيتال}][\ce{H2O}]/([\text{ألدهيد}][\ce{CH3OH}]^2)$: الفائض
الكبير من الميثانول، المربَّع في المقام، يُبقي $Q < K$ حتى لا يبقى إلا قليل من
الألدهيد. أما إزالة الماء فتؤثر في البسط؛ واستعمال الميثانول
مذيبًا أبسط لكنه يحتاج إلى ناتج يسهل فصله.
\end{solution}

\begin{solution}{exo:b1:acetals-protection:12}
لا يستطيع أن يغادر على هيئة ماء إلا مجموعة OH \omterm{def:b1:acetals-protection:hemiacetal}{نصف الأسيتال}، فيتكوّن كاتيون تثبّته
ذرة أكسجين الحلقة المجاورة (\ce{C=O+})؛ أما مجموعات OH الأخرى فكان عليها أن
تغادر من ذرات كربون عادية، فتعطي كاتيونات غير مثبَّتة. ويُضاف الميثانول
إلى الكاتيون المثبَّت: \omterm{def:b1:acetals-protection:hemiacetal}{أسيتال} ميثيلي.
\end{solution}

\begin{solution}{pb:b1:acetals-protection:1}
\textbf{1.} \ce{BrMgCH2CH2CH2CHO}.
\textbf{2.} كربونه الشبيه \omterm{def:b1:electronic-effects:intermediates}{بالأنيون الكربوني} كان سيهاجم ألدهيد جزيء
آخر (أو ألدهيده هو): يُضاف \ce{RMgBr} إلى $\mathrm{R'CHO}$، فيعطي ألكوكسيدًا.
\textbf{3.} الألدهيد، ضد \omterm{def:b1:organomagnesium:organometallic}{كاشف غرينيار}.
\textbf{4.} \omterm{def:b1:acetals-protection:hemiacetal}{الأسيتالات} ثابتة تجاه كواشف غرينيار والقواعد
ويزيلها الحمض المائي.
\textbf{5.} يتكوّن بصورة أتم (ثنائي أول واحد بدل جزيئين من الكحول)
ويسهل حلمأته.
\textbf{6.} $15.09/150.9 = \qty{0.1000}{mol}$.
\textbf{7.} $1.20 \times 0.1000 \times 62.0 = \qty{7.44}{g}$.
\textbf{8.} \ce{BrCH2CH2CH2CHO + HOCH2CH2OH <=> C6H11BrO2 + H2O}
(2-(3-بروموبروبيل)-1,3-ديوكسولان).
\textbf{9.} ينشّط الكربونيل ويتيح فقد الماء في
المرحلة الثانية؛ ويُستعاد.
\textbf{10.} يتقطر التولوين والماء معًا، فيتكاثفان ويسقطان في
المجمِّع؛ فيهبط الماء ويبقى، ويفيض التولوين عائدًا.
\textbf{11.} يحتوي $Q$ على $[\ce{H2O}]$ في بسطه؛ وإزالة الماء تُبقي
$Q < K$، فيتقدم التفاعل.
\textbf{12.} الماء أكثف ولا يمتزج مع التولوين: فيتجمع في
القاع؛ ولا تبلغ ذراعَ العودة إلا الطبقة العليا.
\textbf{13.} حين يتوقف تجمّع الماء، عند الحجم المتوقع.
\textbf{14.} $0.1000 \times 194.9 = \qty{19.5}{g}$.
\textbf{15.} \ce{C6H11BrO2 + Mg -> C6H11BrMgO2} (إيثر جاف).
\textbf{16.} يُضاف الكاشف إلى \ce{(CH3)2CO}: \omterm{def:b1:alcohol-activation:alkoxide}{الألكوكسيد}
\ce{(CH3)2C(O^-)CH2CH2CH2-} على السلسلة المحمية، مع \ce{MgBr+}.
\textbf{17.} كان الحمض سيزيل \omterm{def:b1:acetals-protection:protecting-group}{مجموعة الحماية} قبل الأوان وقد
ينزع الماء من الكحول الثالثي.
\textbf{18.} $0.0700 \times 174.0 = \qty{12.2}{g}$.
\textbf{19.} \babelsublr{\foreignlanguage{arabic}{\omterm{def:b1:acetals-protection:hemiacetal}{الأسيتال}} + \ce{H2O} $\to$ \foreignlanguage{arabic}{الألدهيد + الإيثان-1,2-ثنائي أول}}
(\omterm{def:b1:elementary-steps:catalyst}{حفّاز} حمضي).
\textbf{20.} \omterm{def:b1:alcohol-activation:dehydration}{نزع الماء} منه يحتاج إلى \omterm{def:b1:predominant-reaction:strong-acid}{حمض قوي} وحرارة؛ أما الحمض المائي
الخفيف البارد فيحلمئ \omterm{def:b1:acetals-protection:hemiacetal}{الأسيتال} أسرع بكثير.
\textbf{21.} تُضاف OH الكربون 5 إلى كربون الألدهيد 1: حلقة سداسية
(خمس ذرات كربون وذرة أكسجين واحدة) مع مجموعتي ميثيل على الكربون
المجاور لأكسجين الحلقة ومجموعة OH على الكربون الآخر المجاور له.
\textbf{22.} $0.0700 \times 0.90 = \qty{0.0630}{mol}$، أي \qty{8.19}{g}.
\textbf{23.} $0.0630/0.1000 = \qty{63}{\percent}$ (مع اعتبار \omterm{def:b1:acetals-protection:protecting-group}{الحماية}
تامة).
\textbf{24.} $0.1000 \times 18.0 = \qty{1.80}{g}$.
\textbf{25.} $1.80/0.997 =$ \textbf{\qty{1.8}{mL} من الماء} في المصيدة.
\end{solution}
